\section{Motivation and Scope}

The previous version of this note used a reduced-form New Keynesian system as
the accounting skeleton. That is not the right primitive environment for a
classical BCA exercise. The canonical BCA framework of Chari, Kehoe, and
McGrattan is built from a prototype neoclassical growth model, not from price
rigidity, a Phillips curve, or a monetary policy rule. This version therefore
removes the price-setting block entirely.

The exercise now starts from four canonical wedges:
\begin{enumerate}[label=(\roman*)]
    \item the efficiency wedge, entering the production function;
    \item the government or resource wedge, entering the resource constraint;
    \item the labor wedge, entering the intratemporal labor condition; and
    \item the investment wedge, entering the intertemporal capital Euler
    equation.
\end{enumerate}
The primitive uncertainty shock is stochastic volatility in total factor
productivity. Since TFP itself is the efficiency wedge, TFP uncertainty is first
and foremost a second-moment shock to the efficiency wedge. The question is
whether the second-order risk terms generated by this volatility shock can be
represented as movements in the other canonical BCA wedges, especially the
investment wedge.

The intended bridge to the larger project is:
\begin{align*}
    \text{TFP volatility shock}
    &\Rightarrow
    \text{second-moment movement in the efficiency wedge}
    \\
    &\Rightarrow
    \text{risk corrections in nonlinear equilibrium conditions}
    \\
    &\Rightarrow
    \text{canonical BCA wedges}
    \\
    &\Rightarrow
    \text{BCA-style counterfactuals}.
\end{align*}

This note is still a diagnostic step. It does not yet introduce small-open-
economy trade wedges or UIP wedges, and it does not impose the later empirical
two-factor structure from the OI-SVMVAR. It focuses on one primitive volatility
source so that the mapping from uncertainty to BCA wedges is transparent.
